%% "reactive systems" in everyday life
When we think of computers, we usually picture desktops, laptops, or maybe smartphones. But our
interaction with computer-based systems extends far beyond these familiar devices. From coffee
machines to modern vehicles, many systems around us continuously and automatically react to
\emph{real-time} information. Such systems are naturally referred to as \emph{reactive systems}.

\begin{wrapfigure}{r}{0.5\textwidth}
    \centering
    \includegraphics[width=0.4\textwidth]{reactive_system.png}
    \caption[Reactive system illustration.]%
    {A reactive system illustrated as a coffee machine: the programmer requests coffee, the
        machine produces it, the programmer drinks, and requests again, \textit{ad infinitum}.}
    \label{fig:intro:topic:reactive_systems}
\end{wrapfigure}

Reactive systems are designed to respond immediately to changes in their environment, performing the
reaction loop depicted in \Cref{fig:intro:topic:reactive_systems}. These systems constantly monitor
the flow of inputs and adjust their behavior---whether handling a coffee request or responding to
throttle pedal pressure. The goal is clear: to react promptly and correctly to external events.
%
The level of guarantees expected to achieve this goal varies across reactive systems. Some systems
are considered critical and must provide strong assurances of correct behavior under all conditions,
while others can tolerate occasional failures.
%
For example, if a coffee machine fails to operate one morning, it is merely an inconvenience, not as
serious a hazard as a car blocked in cruise control~\footnote{%
    \href{https://www.ladepeche.fr/article/2013/02/12/1559158-son-regulateur-de-vitesse-bloque-il-franchit-3-peages-a-200-km-h.html}%
    {A real-world example of a blocked cruise control.}%
}.

This thesis addresses the programming of \emph{critical} reactive systems \emph{embedded} in
automotive applications, such as the Adaptive Cruise Control or anti-lock braking. The research was
conducted at \ampere, the software division of the Renault Group.

\section{Domain Specific Languages}
\label{sec:intro:dsl}

Programming critical reactive systems is inherently complex and error-prone. These systems
often handle multiple concurrent processes under strict timing constraints, and must operate
within tight computational limits, such as restricted memory and processing power. In embedded
applications---like those found in automotive systems---these constraints are particularly severe:
the software must be lightweight and efficient to execute on constrained hardware.

\subsection{Low-Level Programming and Its Challenges}
\label{sec:intro:dsl:challenge}

\begin{wrapfigure}{l}{0.5\textwidth}
    \centering
    \includegraphics[width=0.5\textwidth]{languages_level.png}
    \caption[Abstraction level of programming languages.]%
    {A classification of programming languages by abstraction level: lower levels are closer
        to hardware; higher levels offer greater abstractions.}
    \label{sec:intro:dsl:level}
\end{wrapfigure}

To meet these constraints, embedded systems are implemented using \emph{low-level} programming
languages such as \Ci or even assembly.
%
A (subjective) classification of language abstraction levels is illustrated in
\Cref{sec:intro:dsl:level}. In this classification, the lower levels correspond to languages closest
to machine code---those that operate near the hardware and binary encoding. These provide
fine-grained control over resources and allow for highly optimized executables. As we move upward,
we encounter languages that offer increasingly human-friendly abstractions, such as functions, data
types, or even graphical models like block diagrams.
%
Programming languages tend to evolve by building on prior ones, gradually introducing higher-level
abstractions. As a result, boundaries become blurry, and opinions differ on whether \Ci is
best described as low-level or high-level~\footnote{%
    The classification of \Ci as low-level or high-level varies depending on context; it provides
    abstractions above assembly but also allows fine-grained control close to
    hardware (\eg its description in \href{https://en.wikipedia.org/wiki/C_(programming_language)}
    {\textsc{Wikipedia}}).%
}.

However, programming directly in low-level languages is challenging and increases the likelihood of
subtle bugs. Such bugs can have severe consequences. The infamous case of the \emph{Therac-25}
radiation therapy machine exemplifies the risks involved~\cite{journals/computer/LevesonT93}:
several patients received lethal overdoses of radiation due to software errors in a system
programmed directly in assembly language.
%
The use of low-level code made the software difficult to understand, test, and maintain. Critical
failures arose from issues like race conditions---where different parts of the program accessed
shared memory at the same time without proper coordination---and undetected overflows in memory.
Because the system lacked basic safety checks, small timing errors could lead to dangerous
configurations going unnoticed.
%
This tragic example highlights how low-level programming in critical systems can lead to
errors that are hard to detect, with life-threatening consequences.

\subsection{Domain-Specific Languages as a Solution}
\label{sec:intro:dsl:solution}

Given these challenges, there is a strong need for correctness and reliability in critical reactive
systems. Relying solely on low-level programming is often insufficient to guarantee safety.
\emph{Domain-Specific Languages} (\dsl{}s) address this challenge by providing specialized,
high-level abstractions tailored to a specific problem domain.
%
For instance, \href{https://en.wikipedia.org/wiki/SQL}{SQL} is a well-known DSL used for managing
and querying relational databases. Rather than requiring low-level loops or memory access, SQL
allows to express data operations declaratively: for example, the query
\texttt{SELECT name FROM users WHERE active = TRUE;} returns the name of users that are active.

By offering a readable and concise syntax with domain-specific semantics, a \dsl helps developers
specify behavior more clearly. This clarity not only reduces the risk of errors but also improves
maintainability, since a program that is easy to understand is also easier to modify correctly.
Furthermore, a \dsl facilitates automated verification, as analysis tools can be tailored to its
specific constructs. For safety-critical systems, this simplifies certification: if the compiler for
the \dsl is certified, validation efforts can focus on the high-level source code.

\subsection{Dataflow Models for Reactive Systems}
\label{sec:intro:dsl:dataflow}

In the context of \dsl{}s for reactive systems, a common foundation is the \emph{dataflow model},
where systems are described as networks of components that process flows of data over time. This
model is especially well-suited to reactive systems because their inputs---such as sensor readings,
user interactions, or communication messages---naturally arrive as flows of information, one value
at a time, in a temporal sequence. The dataflow model offers an expressive way to describe reactive
behavior by emphasizing how data moves through the system, rather than dictating the precise
sequence of operations for the computer to perform.

\paragraph{Kahn Process Network}
\label{sec:intro:dsl:dataflow:kpn}

One well-known formalization of dataflow is the \emph{Kahn Process Network} (KPN)
model~\cite{DBLP:conf/ifip/Kahn74}, in which independent concurrent processes communicate through
unbounded FIFO (First-In, First-Out) channels. Channel reads are blocking (a process waits if no
data is available), while writes are non-blocking thanks to the assumption of infinite channel
capacity. Processes cannot test whether a channel is empty, enforcing strict data dependency. If all
processes are deterministic, then the entire network is deterministic as well: output dataflows
depend only on input dataflows, not on the execution order or communication timing.

A major drawback of KPNs is the lack of guarantees on deadlock-freedom (processes may wait
indefinitely for input from each other) and bounded memory usage (unbounded channels may grow
without limit if data production outpaces consumption, making memory usage unmanageable in embedded
systems).

\paragraph{Synchronous Dataflow Models}
\label{sec:intro:dsl:dataflow:sdf}

To address these drawbacks, one approach is to restrict models to a subset of KPNs that can be
statically analyzed for memory boundedness. This is the case with \emph{Synchronous Dataflow} (SDF)
models~\cite{DBLP:journals/tc/LeeM87}, which allow compile-time scheduling of process execution and
bounded buffer sizes. These properties make SDF particularly well-suited for digital signal
processing (DSP) applications. As a result, SDF has been adopted in industrial DSP tools such as
COSSAP and the Signal Processing WorkSystem~\cite{books/kap/BattacharyyaLM96}.

\paragraph{Synchronous Dataflow Languages}
\label{sec:intro:dsl:dataflow:synchronous_languages}

\Lustre~\cite{DBLP:conf/popl/CaspiPHP87} and its industrial implementation
\Scade~\cite{DBLP:conf/tase/ColacoPP17} were developed to program critical embedded systems using
a declarative style. In these \emph{synchronous dataflow languages}, KPNs are restricted
to eliminate communication buffers entirely: each process is associated with a clock that defines
when computations occur. Flows can be filtered and combined using clock operators, and all
timing dependencies are made explicit in the program structure. The \emph{synchronous hypothesis} is
assumed: outputs are produced instantaneously with respect to inputs, simplifying concurrency
modeling. A synchronous program reacts at discrete logical instants governed by its \emph
{base clock}, and compilation ensures bounded memory usage and bounded reaction time. The latter can
be measured as the \emph{Worst Case Execution Time} (\wcet), which is checked to be smaller than the
base clock period to verify the synchronous hypothesis and guarantee real-time correctness.

Other synchronous languages include \Signal~\cite{DBLP:journals/scp/BenvenisteGJ91} and
\Esterel~\cite{DBLP:journals/scp/BerryG92}. While \Lustre is strictly functional, \Signal is
relational and emphasizes a fine-grained control over timing by allowing multiple clocks and
explicit equations over their relations. This clock calculus allows for modeling distributed
and multi-rate systems, but requires solving logical constraints at compile time, making
compilation more complex. In contrast, \Esterel follows an imperative style and offers much richer
expressivity, including complex control structures such as preemption and exception handling.
Its semantics are compiled into boolean circuits~\cite{DBLP:books/daglib/0023221}, making it
suitable for hardware synthesis and time-critical control logic.

\paragraph{Functional Reactive Programming}
\label{sec:intro:dsl:dataflow:frp}

Another family of dataflow-based \dsl{}s is \emph{Functional Reactive Programming}
(\frp)~\cite{DBLP:conf/icfp/ElliottH97}, initially introduced to describe interactive animations and
graphical user interfaces (GUIs) in a declarative and composable way. In \frp, programs manipulate
time-varying values, such as user input or sensor data, using high-level functional abstractions.
These values are represented as \emph{signals} (continuous or discrete over time) and \emph{events}
(discrete occurrences), allowing the programmer to define reactive behavior as pure functions over
dataflows. Unlike synchronous languages like \Lustre, which assume a single base clock and bounded
reaction times, \frp embraces a more flexible and often asynchronous notion of time. It prioritizes
expressivity and composability over predictability and timing guarantees.
%
Because of this focus, \frp languages are typically used in non-critical domains, such as user
interfaces, multimedia systems, or robotics. Popular \frp implementations include
\Yampa~\cite{DBLP:conf/haskell/CourtneyNP03} in \Haskell, \Elm~\cite{DBLP:conf/pldi/CzaplickiC13}
in \JavaScript, and \ReactiveX libraries~\cite{reactivex} for multiple languages. Although \frp
shares the dataflow spirit with synchronous languages, its execution model is generally based on
dynamic data propagation and lazy evaluation, rather than statically scheduled computations. This
makes formal timing analysis and resource prediction more difficult, but enables elegant and highly
modular code for reactive applications.

Early \frp systems lacked formal guarantees, particularly regarding \emph{causality}---ensuring that
outputs depend only on past and present inputs---and bounded use of memory and time. These
limitations made them difficult to apply in real-time or resource-constrained environments. To
address this, several \frp variants have been proposed with stronger guarantees. Real-Time
\frp~\cite{DBLP:conf/icfp/WanTH01} introduces restrictions to ensure that computations consume
bounded memory and execute within a predictable time. Event-Driven
\frp~\cite{DBLP:conf/padl/WanTH02} addresses causality by making all updates explicitly triggered by
discrete events, which ensures deterministic and well-defined signal propagation. Other variants,
such as Modal \frp, use a modal type theory to formally prevent issues like memory leaks, as
demonstrated by \textsc{Rattus}~\cite{DBLP:journals/pacmpl/BahrGM19}, which enforces a clock-based
execution model, and its asynchronous extension
\textsc{Async Rattus}~\cite{DBLP:journals/pacmpl/BahrM23}. These
refinements make \frp more suitable for embedded and reactive systems while preserving its
declarative and compositional nature.

\paragraph{Execution Strategies in Reactive Systems}
\label{sec:intro:dsl:dataflow:push_pull}

Execution in reactive systems can follow two main strategies: \emph{push}-based or
\emph{pull}-based~\cite{DBLP:conf/haskell/Elliott09}. In push-based systems, also called
\emph{event-driven}, updates are propagated eagerly whenever inputs change. This makes them
responsive to external stimuli, but can be inefficient when inputs change frequently and the
results are not always needed. In contrast, pull-based or \emph{demand-driven} systems compute
outputs only when they are explicitly requested. This avoids unnecessary work and provides
predictable computation, but may introduce latency due to on-demand evaluation (if the outputs had
been updated eagerly, they would already be available) and redundant computations (if inputs have
not changed).
%
This distinction influences how data dependencies propagate and how resources are managed. Most
synchronous language compilers adopt a pull-based approach: computations occur at each tick of the
system's base clock, independently of input changes. This supports static scheduling and
predictable resource usage. In contrast, many \frp implementations support push-based or hybrid
strategies~\cite{reactivex,DBLP:conf/pldi/CzaplickiC13,DBLP:conf/ppdp/MandelPP15}. Push-based
execution suits input-driven systems like GUIs or sensor networks, while pull-based evaluation
enables on-demand computation. Hybrid models aim to combine both advantages by avoiding unnecessary
updates and improving efficiency~\cite{DBLP:conf/haskell/Elliott09}.

\subsection{Summary}
\label{sec:intro:dsl:summary}

This section showed how critical reactive systems require specialized languages to handle
complexity and strict constraints. Domain-specific languages provide high-level abstractions
tailored for these needs, often based on dataflow models and synchronous execution semantics.
With foundations in place, we now examine how these languages are applied in the automotive
industry, where safety and timing demands are paramount.

\section{Programming Cars}
\label{sec:intro:cars}

Modern vehicles rely on a large number of embedded software components to implement their
functionalities, from engine control and braking systems to infotainment and driver assistance.
These components are inherently reactive: they continuously interact with the physical world,
respond to sensor inputs, issue actuator commands, while operating concurrently under strict timing
and resource constraints.

\subsection{Standards and Architectures}
\label{sec:intro:cars:standards}

To guarantee functional safety, most automotive manufacturers adhere to the \iso standard,
which provides a systematic framework for identifying and mitigating risks in road vehicles.
It introduces the notion of \emph{Automotive Safety Integrity Levels} (\Asil), ranging from A
(least critical) to D (most critical), based on the severity, exposure, and controllability
of potential hazards. For example, direction indicators may be assigned \Asil~A, while an
engine controller is typically classified as \Asil~D. Higher \Asil levels impose stricter
requirements on testing, traceability, determinism, and tool qualification.

The \ausar (\emph{AUTomotive Open System ARchitecture}) standard was introduced by major
automotive manufacturers and suppliers to define a common interface between software and hardware
components. It provides a modular architecture for \emph{Electronic Control Units} (\ecus)---small
embedded computers responsible for specific vehicle functions such as engine management, braking,
or lighting. \ausar structures software into standardized layers to promote reuse and hardware
independence. At its core is the \emph{Runtime Environment} (RTE), a middleware layer that
abstracts communication between application components and the lower-level infrastructure. Above
the RTE are the \emph{Application Software Components}, which implement functional logic, while
below it lies the \emph{Basic Software}, responsible for hardware interaction, communication,
and diagnostics. This layered design enables independent development, testing, and integration of
software modules, improving interoperability across suppliers and vehicle platforms.

In typical deployments, automotive software is distributed across multiple \ecus, which interact
\via communication buses such as CAN, LIN, or Ethernet. Each \ecu operates under a periodic
schedule: it samples its inputs at regular intervals, computes outputs, and sends results to other
components. This execution model aligns well with synchronous programming languages.

\subsection{Model-Based Development}
\label{sec:intro:cars:mbd}

To manage complexity and ensure correctness, automotive software is often developed using
\emph{Model-Based Development} techniques. In this approach, engineers describe component behavior
using graphical or textual models in high-level \dsl{}s. These models are used to simulate behavior,
verify safety properties, and generate \Ci code automatically.

One of the most widely used tools in industry is \emph{\Simulink}~\cite{simulink}, a graphical
dataflow-based environment for modeling systems. \Simulink allows engineers to connect blocks
representing signals, operators, or state machines. Its integration with \Stateflow enables control
flow modeling, while code generation toolchains produce production-grade \Ci code. \Simulink models
can be configured to comply with \iso and integrate with \ausar through dedicated toolboxes.

While \Simulink shares similarities with synchronous dataflow languages (such as its use of discrete
time steps and block diagrams) it does not follow a formally defined synchronous semantics. Instead,
its behavior depends on block sample times, solver settings, and execution order, which must be
carefully configured to ensure predictable and deterministic execution.

\subsection{Integration of Synchronous \dsl{}s}
\label{sec:intro:cars:synchronous}

For safety-critical components, synchronous languages such as \Lustre and its industrial
derivative \Scade offer strong formal guarantees. These \dsl{}s ensure deterministic behavior,
bounded memory usage, and analyzable timing, which are essential for high-integrity real-time
systems. Their compilation process produces statically scheduled code with a guaranteed
\emph{Worst-Case Execution Time} (\wcet), enabling formal verification against timing constraints
imposed by \iso and \ausar models.

\Scade has been widely adopted in avionics and nuclear industries, where the highest certification
levels (\eg \textsc{DO-178C} Level A for aerospace) are required. It is also qualified under the
\iso standard up to \Asil~D, and although its use in the automotive domain is more limited,
it is applicable to the most critical software components. In practice, however, automotive systems
more commonly rely on \Simulink or other modeling tools, which offer greater flexibility and are
tightly integrated with the \ausar ecosystem. These tools often trade some formal rigor for
ease of use and broader industry support, making them suitable for less critical functionality such
as driver assistance or infotainment.

\subsection{Summary}
\label{sec:intro:cars:summary}

Developing embedded software for modern vehicles requires balancing safety, reliability, and
engineering complexity. Standards such as \iso and \ausar establish structured foundations for
designing, validating, and integrating automotive software components. Model-based development tools
like \Simulink facilitate rapid prototyping and integration within the \ausar ecosystem, while
formally defined synchronous \dsl{}s such as \Lustre and \Scade provide stronger guarantees for
the most critical subsystems. Together, these technologies support the scalable and safe development
of increasingly complex automotive software.

\section{Architectural Shift to Software-Defined Vehicles}
\label{sec:intro:gap}

Advanced Driver Assistance Systems, computer vision, sensor fusion, and AI-based features are
becoming central to modern vehicle functionality. These capabilities significantly increase the
complexity and computational demands of in-vehicle software. To address these challenges, \ampere is
moving away from traditional distributed electronic architectures and adopting a centralized
\emph{Software-Defined Vehicle} (\sdv) architecture.

\begin{figure}[ht!]
    \centering
    \begin{subfigure}[t]{0.4\textwidth}
        \centering
        \includegraphics[height=\linewidth]{ecu_archi.png}
        \caption[Legacy architecture of automotive software.]%
        {Legacy architecture, featuring numerous specialized \ecus (grey boxes).}
        \label{fig:intro:gap:archi:ecus}
    \end{subfigure}
    \quad
    \begin{subfigure}[t]{0.4\textwidth}
        \centering
        \includegraphics[height=\linewidth]{central_archi.png}
        \caption[\emph{Software Defined Vehicle} architecture.]%
        {\sdv architecture with centralized computers (grey box) connected
            to sensors and actuators (grey circles).}
        \label{fig:intro:gap:archi:sdv}
    \end{subfigure}
    \caption[Legacy \versus Software Defined Vehicle architecture.]%
    {A comparison of the legacy \ecu architecture and the modern \sdv architecture.}
    \label{fig:intro:gap:archi}
\end{figure}

The legacy architecture, illustrated in \Cref{fig:intro:gap:archi:ecus}, relies on a distributed
system of specialized Electronic Control Units (\ecus). In this model, introducing new functionality
often required integrating a dedicated \ecu. This approach made it difficult to add
computation-intensive features due to physical constraints, such as finding space for new hardware,
managing complex wiring, and scheduling inter-\ecu communication. Moreover, this design lacks
support for software updates, a critical limitation for rapidly evolving AI-based features.

In contrast, the \sdv architecture centralizes core functionalities on high-performance computers,
as shown in \Cref{fig:intro:gap:archi:sdv}. In this model, functionalities are implemented as
software that can be updated continuously, often remotely. This architectural shift away from
numerous individual \ecus reduces physical complexity, cabling, and overall vehicle weight. The
resulting design is more efficient, since computation is consolidated on powerful hardware, and more
maintainable, as updates no longer require physical component replacement. This approach provides
lower manufacturing costs, reduced energy consumption, and a foundation for continuous software
innovation.

However, this shift exposes limitations in the traditional embedded software development toolchain.
A common approach in the automotive domain involves model-based development with tools such as
\Simulink and \Stateflow, which typically generate standard-compliant \Ci code based on a periodic
execution model. While effective for safety-critical systems in traditional architectures, this
approach presents several challenges in the context of \sdv.
%
\begin{description}
    \item[Inefficient execution:] Periodic scheduling allocates fixed time slices for task
          execution, regardless of whether new data is available. This leads to redundant
          computations and unnecessary data transfers, wasting CPU cycles and saturating
          communication buses.

    \item[Manual optimization burden:] The generated \Ci code typically requires hand-tuning to meet
          the timing and resource constraints. These manual adjustments jeopardize maintainability
          and scalability, especially as systems grow in complexity.

    \item[Costly safety verification:] Verifying safety properties at the \Ci code level is
          time-consuming and costly. The lack of high-level semantic information in generated code
          complicates formal analyses, particularly for large systems.

    \item[Limited expressiveness:] Traditional tools offer limited support for modeling event-driven
          behavior, which is increasingly common in \sdv software. Capturing such behavior with
          periodic models often results in complex implementations.
\end{description}

\subsection{The Interest in \Rust}
\label{sec:intro:gap:rust}

The shift to centralized computing platforms significantly affects the software development
lifecycle. Integration becomes more complex as functionalities that were once isolated on separate
\ecus now share common hardware resources. This increases the need for advanced resource
management strategies and strong fault isolation to ensure safety and system robustness.

To address these challenges, there is growing interest in general-purpose programming languages that
offer strong memory safety, concurrency control, and low-level performance. \ampere, along with
other automotive industry players~\cite{ferrocene,veecle}, is progressively adopting the \Rust
programming language. This choice represents a strategic shift away from traditional \Ci code. While
performant, \Ci provides no memory safety guarantees. Ensuring the correctness of \Ci code requires
extensive verification through tool qualification, static analysis, or formal proofs---a process
that is both tedious and fragile.

In contrast, \Rust provides safety \emph{by construction}, embedding formal-methods principles
directly into the language design. Its compiler enforces memory and thread safety at compile time
through a powerful system of \emph{ownership}, \emph{borrowing}, and \emph{lifetimes}. The
\emph{ownership} model ensures that every value has a unique owner, preventing memory leaks and
premature deallocations. The \emph{borrow checker} statically guarantees that references never
outlive the data they point to and enforces a key rule: data can have either multiple immutable
references or one mutable reference, but not both. This rule eliminates entire classes of
bugs---such as memory corruption, buffer overflows, and data races---without the runtime overhead of
a garbage collector.

For low-level operations that are inherently unsafe, such as direct hardware interaction, \Rust
provides the \texttt{unsafe} keyword. This feature allows programmers to bypass certain compiler
checks in a delimited block, explicitly marking where the responsibility for upholding safety
invariants falls on them. The \RustBelt project~\cite{DBLP:journals/pacmpl/0002JKD18} provides a
formal foundation for verifying such code, enabling the development of proven-safe abstractions over
low-level primitives.

\subsection{Bus Overload at \ampere}
\label{sec:intro:gap:bus_overload}

The transition to an \sdv architecture, with its reliance on high-bandwidth buses like Automotive
Ethernet, introduces challenges not present in legacy architectures. A primary issue is
\emph{protocol overhead}: the extra data required to format, address, and transmit a message, beyond
the useful data---called \emph{payload}.

The CAN bus, designed for automotive use, is extremely lightweight. A CAN message contains a short
payload (up to 8 bytes, 64 bytes with CAN-FD), a message identifier, and a few bits for control and
error-checking. The protocol uses a clever mechanism called non-destructive bitwise arbitration to
manage message priority without a central coordinator. This is possible due to a physical layer
convention: a \texttt{0} bit (dominant) actively drives the bus to a specific voltage, while a
\texttt{1} bit (recessive) is a passive, high-impedance state. If a controller, while sending a
recessive bit, detects a dominant bit on the bus, it knows a higher-priority message is being
transmitted, immediately stops, and becomes a receiver. Since the dominant bit is logically
\texttt{0}, the message with the numerically lowest identifier always wins arbitration without being
corrupted, guaranteeing that critical data (\eg from the brake controller) is never delayed by less
important traffic (\eg from the radio). This makes CAN highly efficient and deterministic
for small, frequent messages.

In contrast, Automotive Ethernet typically uses protocol stacks like \texttt{TCP/IP}, designed for
general-purpose networking. This layered approach introduces significant overhead. To send data, the
system must construct a packet by encapsulating the payload within multiple layers, each adding its
own header:
\begin{itemize}
    \item The \emph{Ethernet frame}, which uses hardware (MAC) addresses to deliver data between
          devices on the same local network segment.
    \item The \emph{IP packet}, which uses logical (IP) addresses to route data across different
          interconnected networks within the vehicle.
    \item The \emph{TCP or UDP segment}, which uses port numbers to direct the data to the correct
          application (\eg the infotainment system or a sensor fusion) on the destination device.
\end{itemize}
As a result, sending just a few bytes of data requires a packet with at least 40-50 bytes of
overhead. This high cost manifested as a critical problem during an experiment at \ampere.


% \begin{table}
%     \newcommand{\rot}[1]{\rotatebox{60}{#1}}
%     \begin{subtable}{0.45\textwidth}
%         \centering
%         \begin{tabular}{rcccccc}
%             \rot{ECU Name}  & \rot{Received Rx Events/s} & \rot{Consumed Rx Events/s} & \rot{Consumed/Received} & \rot{Rx Packets/s} & \rot{Tx Events/s} & \rot{Tx Packets/s} \\\hline
%             \texttt{PCU\_1} & 3123                       & 2259                       & 72\%                    & 3123               & 2451              & 4831               \\
%             \texttt{PCU\_2} & 3052                       & 208                        & 7\%                     & 3052               & 20                & 20                 \\
%             \texttt{PCU\_3} & 785                        & 785                        & 100\%                   & 785                & 2325              & 2385               \\
%             \texttt{PCU\_4} & 805                        & 25                         & 3\%                     & 805                & 20                & 20                 \\
%             \texttt{PCU\_5} & 2849                       & 1915                       & 67\%                    & 2849               & 0                 & 0                  \\
%             \texttt{PCU\_6} & 2849                       & 215                        & 8\%                     & 2849               & 20                & 20                 \\
%             \texttt{PCU\_7} & 2849                       & 835                        & 29\%                    & 2849               & 240               & 9330               \\
%             \texttt{PCU\_8} & 4330                       & 4030                       & 93\%                    & 17620              & 20                & 20                 \\
%             \texttt{PCU\_9} & 2849                       & 709                        & 25\%                    & 2849               & 100               & 2280               \\
%             \texttt{CDC\_1} & 815                        & 40                         & 5\%                     & 815                & 0                 & 0                  \\
%             \texttt{CDC\_2} & 1995                       & 1740                       & 87\%                    & 2055               & 0                 & 0                  \\
%             \texttt{CDC\_3} & 990                        & 990                        & 100\%                   & 990                & 0                 & 0                  \\
%             \texttt{CDC\_4} & 935                        & 675                        & 72\%                    & 935                & 0                 & 0                  \\
%             \texttt{PIU\_1} & 2849                       & 264                        & 9\%                     & 2849               & 104               & 104                \\
%             \texttt{PIU\_2} & 3376                       & 2926                       & 87\%                    & 3376               & 2005              & 2105               \\
%             \texttt{PIU\_3} & 2979                       & 140                        & 5\%                     & 2979               & 30                & 30                 \\\hline
%         \end{tabular}
%     \end{subtable}
%     \hfill
%     \begin{subtable}{0.45\textwidth}
%         \centering
%         \begin{tabular}{rccccc}
%             \rot{ECU Name}  & \rot{Received Rx Events/s} & \rot{Consumed/Received} & \rot{Rx Packets/s} & \rot{Tx Events/s} & \rot{Tx Packets/s} \\\hline
%             \texttt{PCU\_1} & 165                        & 100\%                   & 165                & 3498              & 3498               \\
%             \texttt{PCU\_2} & 20                         & 100\%                   & 20                 & 10                & 10                 \\
%             \texttt{PCU\_3} & 19                         & 100\%                   & 19                 & 1740              & 1740               \\
%             \texttt{PCU\_4} & 2                          & 100\%                   & 2                  & 10                & 10                 \\
%             \texttt{PCU\_5} & 725                        & 100\%                   & 725                & 19                & 19                 \\
%             \texttt{PCU\_6} & 0                          & 100\%                   & 0                  & 10                & 10                 \\
%             \texttt{PCU\_7} & 100                        & 100\%                   & 100                & 211               & 3331               \\
%             \texttt{PCU\_8} & 4744                       & 85\%                    & 8824               & 229               & 229                \\
%             \texttt{PCU\_9} & 0                          & 100\%                   & 0                  & 101               & 1061               \\
%             \texttt{CDC\_1} & 303                        & 13\%                    & 303                & 20                & 20                 \\
%             \texttt{CDC\_2} & 173                        & 100\%                   & 173                & 3                 & 3                  \\
%             \texttt{CDC\_3} & 10                         & 100\%                   & 10                 & 0                 & 0                  \\
%             \texttt{CDC\_4} & 400                        & 100\%                   & 400                & 18                & 18                 \\
%             \texttt{PIU\_1} & 342                        & 77\%                    & 342                & 884               & 884                \\
%             \texttt{PIU\_2} & 1528                       & 100\%                   & 1528               & 3615              & 3615               \\
%             \texttt{PIU\_3} & 845                        & 17\%                    & 845                & 626               & 626                \\\hline
%         \end{tabular}
%     \end{subtable}
% \end{table}

\begin{figure}[ht!]
    \centering
    \begin{subfigure}{0.51\textwidth}
        \centering
        \includegraphics[height=0.24\textheight]{overload/bus_2023.png}
        \caption{Bus traffic with periodic transmission (experimented in 2023).}
        \label{fig:intro:gap:bus_overload:exp:2023}
    \end{subfigure}
    \hfill
    \begin{subfigure}{0.45\textwidth}
        \centering
        \includegraphics[height=0.24\textheight]{overload/bus_reduction_2025.png}
        \caption{Traffic reduction with on-change transmission (experimented in 2025).}
        \label{fig:intro:gap:bus_overload:exp:2025}
    \end{subfigure}
    \caption[Comparison of bus load on a prototype vehicle board at \ampere.]%
    {Comparison of bus load on a prototype vehicle board at \ampere, showing the
        inefficiency of periodic data transmission and the effectiveness of an on-change model.}
    \label{fig:intro:gap:bus_overload:exp}
\end{figure}

During full-scale tests on a vehicle board prototype, engineers observed that the main communication
bus was becoming overloaded. The system, following a traditional periodic model, was configured to
transmit every signal at a fixed 10\,ms interval. This approach, inherited from the predictable
world of CAN, is proved unsustainable. A detailed analysis, summarized in
\Cref{fig:intro:gap:bus_overload:exp:2023}, revealed a critical inefficiency. For some elements, as
few as 3\% of received messages contained information that was effectively used. The remaining 97\%
were redundant, and the high traffic volume led to dropped packets, consuming valuable network
bandwidth and CPU cycles while causing data loss.

To address this, engineers shifted to a hybrid \emph{on-change} communication paradigm. This model
was implemented by applying a filtering algorithm on the receiving end. The receiver's application
is only notified when the value of the data actually changes. To meet safety requirements, this
model is augmented with two temporal constraints:
\begin{itemize}
    \item A \emph{debounce time}, which sets a maximum notification frequency to prevent the
          application from being flooded by rapid signal fluctuations.
    \item A \emph{refresh time}, which enforces a minimum transmission frequency, acting as a
          heartbeat to confirm that a component is still live even if its data has not changed.
\end{itemize}
As shown in \Cref{fig:intro:gap:bus_overload:exp:2025}, this strategy is proved highly effective. By
filtering redundant data at the receiver, engineers reduced the number of events processed by the
application by an average of 75\%, with projections indicating a corresponding 50\% reduction in CPU
load from packet processing.

The experiment thus demonstrated that a purely periodic model is ill-suited for the dynamic,
data-intensive workloads of \sdv architectures, highlighting the need for an event-driven execution
model to manage resources efficiently.

\subsection{The Choice for Asynchronous \Rust}
\label{sec:intro:gap:async_rust}

As demonstrated by the bus overload at \ampere, the traditional periodic execution model is
ill-suited for \sdv architecture. This model wastes CPU cycles and network bandwidth by executing
tasks at fixed intervals, regardless of need. This inefficiency becomes a critical bottleneck when
integrating data-intensive but non-critical (non \Asil~D) functions, such as infotainment and ADAS.

Asynchronous programming presents an attractive solution. By adopting an event-driven model, tasks
can give-back control of the processor when they are idle (\eg waiting for inputs), freeing the CPU
to handle other tasks. This approach is a good fit for the dynamic nature of modern automotive
software, improving overall resource management. Choosing to implement this model in \Rust adds a
crucial benefit: the language memory safety guarantees help isolate components, preventing
less-critical functions from interfering with safety-critical ones. For these reasons, \ampere
selected the \texttt{tokio} runtime~\cite{tokio} to build its next-generation software, leveraging a
mature, high-performance foundation for managing concurrent, event-driven tasks.

However, this decision introduces downsides. First, asynchronous programming is known to be complex
and error-prone. Reasoning about program flow becomes more difficult, and it introduces new classes
of potential bugs, such as deadlocks or incorrect state management across await points. Second, the
approach also creates a strong dependency on the runtime; the correctness and performance of the
system hinge on the implementation of \texttt{tokio}, which uses internal \texttt{unsafe} code for
its operation. Finally, runtimes like \texttt{tokio} are optimized for high throughput, not for the
predictable, low-latency execution required for hard real-time tasks. This makes them unsuitable for
the most critical (\Asil~D) components, creating a fundamental tension between efficiency and
temporal predictability in a mixed-criticality system.

\subsection{The Abstraction Gap}
\label{sec:intro:gap:abstraction}

Despite its strengths, the low-level nature and unfamiliar programming model of \Rust present a
steep learning curve, especially for engineers accustomed to model-based design. Engineers must
adapt to a paradigm that emphasizes fine-grained memory control, ownership semantics, and explicit
concurrency. While the strong compile-time guarantees of \Rust improve memory and thread safety,
they can also make development more difficult, particularly in early stages of adoption. Although
some abstractions are emerging, there remains a lack of high-level modeling tools tailored to
control systems and real-time dataflow. This gap complicates the transition from traditional
development environments.
%
Reactive \Rust libraries such as \texttt{tokio}\footnote{While best known for its runtime,
    the \texttt{tokio} project also provides a rich library of asynchronous primitives.} and
\texttt{futures}~\cite{tokio,futures} provide powerful concurrency primitives, but they operate at a
low abstraction level and are not intended for control system modeling.
%
\textsc{Veecle}~\cite{veecle} provides a modular runtime and high-level \textsc{API}s designed for
\sdv architectures, abstracting middleware and enabling asynchronous execution.
However, it primarily focuses on runtime infrastructure rather than providing modeling abstractions
familiar to control engineers.
%
\textsc{Pictorus}~\cite{pictorus}, in contrast, offers a modeling environment with block diagrams
that generates \Rust code and supports real-time collaboration. Yet, the scheduling model
remains sequential, lacking native support for the event-driven behaviors essential for
modern \sdv systems.

These gaps motivate the development of new modeling tools that bring the rigor and usability of
traditional approaches to the modern \sdv context. By bridging the abstraction gap, such tools can
facilitate a smoother transition to \Rust-based development and support the unique requirements of
centralized automotive computing platforms.

\section{Contribution}
\label{sec:intro:contrib}

To address these challenges, we introduce \GRust, a \emph{Domain-Specific Language} (\dsl) tailored
for automotive system modeling in the context of \sdv architectures. \GRust bridges the gap between
the high-level abstractions of model-based design and the safety, performance, and concurrency
features of \Rust. It provides the following:
%
\begin{itemize}
    \item A declarative syntax based on the \emph{Synchronous} paradigm, underlying tools like
          \Simulink and \Scade, and academic languages such as \Lustre. This model is familiar to
          control engineers and ensures deterministic, step-based computation.
    \item Constructs inspired by \emph{Functional Reactive Programming} (\frp), particularly the
          \textsc{ReactiveX} library, enabling the explicit modeling of asynchronous and
          event-driven behavior with clear data and temporal dependencies.
    \item A compiler that translates \GRust into safe and asynchronous \Rust code, suitable for
          integration with modern, performance-critical \Rust-based \sdv software stacks. This
          automated generation is critical, as writing asynchronous code by hand is notoriously
          complex and error-prone. The compiler shields developers from these low-level details,
          ensuring the generated code is correct by construction, thereby reducing the risk of
          subtle bugs like deadlocks or incorrect state management.
    \item A formal foundation comprising both an \emph{operational semantics}, which models the
          step-by-step execution, and a \emph{denotational semantics}, which provides a mathematical
          interpretation of programs. These semantics, proved to be equivalent, are used to prove
          key safety properties, such as productivity and determinism.
\end{itemize}

Traditional synchronous languages offer deterministic semantics and compile-time
guarantees well suited for embedded systems, and have seen industrial adoption through tools like
\Scade. However, features such as fine-grained clock control often demand deep familiarity with
their formal semantics, limiting their accessibility. These languages also lack native support for
event-driven behavior and do not integrate with the \Rust ecosystem, reducing their applicability to
\sdv platforms.

Conversely, \frp enables high-level modeling of asynchronous dataflows. Research efforts have made
\frp more predictable and resource-efficient by incorporating ideas from synchronous languages,
such as static scheduling~\cite{DBLP:conf/haskell/NilssonCP02,DBLP:conf/ppdp/ChupinN19} and
memory-safe stream processing~\cite{DBLP:conf/icfp/Krishnaswami13}. \GRust adopts the reverse
strategy: it starts from the deterministic structure of synchronous models familiar to control
engineers, and augments it with reactive constructs inspired by \textsc{ReactiveX} to support
event-driven behavior.

\subsection{Short Presentation of \GRust}
\label{sec:intro:contrib:grust}

The \GRust language consists of two sub-languages: the service-modeling kernel \GRFRP and the
component-modeling kernel \GRSync.

\paragraph{Services in \GRFRP}
\label{sec:intro:contrib:grust:grfrp}

\GRFRP defines the system's high-level structure through \emph{services}, which compute dataflows
and interface with the rest of the system \via declared \emph{imports} and \emph{exports}. Inspired
by the \frp paradigm, \GRFRP allows developers to specify flow computations declaratively.
%
In \GRFRP, flows are classified as either \emph{signals}, which represent values that persist and may
change over time, or \emph{events}, which represent punctual values emitted at discrete moments.
These flows are asynchronous: services do not receive all inputs simultaneously, but rather in
independent and possibly irregular drops---each representing a new input value or event occurrence.
%
Each service declares an \emph{interface} by importing flows produced by external entities (\eg
sensors or other services) and exporting computed flows for use elsewhere in the system (\eg
actuators or other services). The behavior of exported flows is defined using operations on imported
ones.
%
Flow operations use either built-in combinators provided by the \GReact library or custom components
defined in the \GRSync sub-language. This separation of concerns supports modular design and reuse.

In essence, a \GRFRP service acts as a bridge between the system's asynchronous dataflows and the
computational logic encapsulated in reusable components.

\paragraph{Components in \GRSync}
\label{sec:intro:contrib:grust:grsync}

\GRSync is a specification language for defining user-defined flow operators, known as
\emph{components}. A component describes the computation of output flows at each logical instant,
based on the current and past values of its input flows.
%
From a conceptual standpoint, \GRSync components are similar to those in
\Lustre~\cite{DBLP:journals/pieee/HalbwachsCRP91}, but without explicit user-defined clocks. For
engineers in industry, \GRSync resembles \Simulink block diagrams~\cite{simulink},
where graphical blocks define operations and wires represent data dependencies. In hardware terms,
components can be thought of as logical chips, with inputs and outputs corresponding to the chip's
pins.
%
Components are reusable building blocks that can be instantiated either within services or other
components. Each instance maintains its own internal state and executes independently, ensuring
isolated behavior. This modularity allows different input flows to be processed concurrently,
producing separate and predictable outputs.

In essence, a \GRSync component captures a deterministic, synchronous computation unit that can be
composed and reused to implement the reactive logic of \GRFRP services.

\paragraph{Globally Asynchronous, Locally Synchronous}
\label{sec:intro:contrib:grust:gals}

The synchronicity of components in \GRSync is not in conflict with the asynchrony of services in
\GRFRP; rather, they reflect a well-established architectural principle known as \emph{Globally
    Asynchronous, Locally Synchronous} (GALS)~\cite{DBLP:phd/us/Chapiro85}. In GALS systems,
individual components operate under their own local synchronous clocks, while communication between
them occurs asynchronously. This design mirrors the physical reality of embedded systems, where
different process operate at varying rates and must communicate reliably despite temporal
mismatches.
%
In \GRust, services act as the asynchronous boundary between the system's external dataflows and the
locally synchronous logic encapsulated in components. They provide a unified abstraction for
handling asynchronous inputs (\eg from sensors or other services) and exposing them to
deterministic, step-based computation models. This integration supports formal analysis, predictable
execution, and safe embedding in event-driven architectures.

\subsection{Prototype Compiler}
\label{sec:intro:contrib:proto}

The documentation for the language is available at:
\begin{center}
    \URLgrbook .
\end{center}
%
The source code of the \GRust compiler can be found at:
\begin{center}
    \URLgrust .
\end{center}
It is implemented as a \Rust procedural macro, enabling seamless integration with the \Rust
environment.
%
% Additionally, the compiler is published as a \Rust library on \texttt{crates.io}:
% \begin{center}
%     \url{https://crates.io/crates/grust} .
% \end{center}

\subsection{Formal Semantics and Safety Guarantees}
\label{sec:intro:contrib:semantics}

Our contribution includes a formalization of \GRust through two complementary semantic models
(operational and denotational), along with theorems that ensure correctness, determinism, and
semantic consistency. Together, these results form a sound theoretical foundation that reinforces
\GRust's suitability for high-assurance embedded software.

\paragraph{Operational Semantics}
\label{sec:intro:contrib:semantics:opsem}

We provide an operational semantics that models the push-based execution of \GRust services, which
react to the arrival of inputs from the system's communication bus. This semantics describes how
changes in flows are propagated within services and their constituent components.

\paragraph{Denotational Semantics}
\label{sec:intro:contrib:semantics:denot}

We complement this with a denotational semantics in which flows are interpreted as infinite
sequences of timestamped values, called \emph{timestamped streams}. For signals, the stream contains
all value changes; for events, it captures the individual occurrences. Flow operators are
interpreted as stream functions. This model is inspired by the theory of synchronous
programming~\cite{DBLP:journals/tecs/BourkePP23} and closely resembles the Tagged Signal
Model~\cite{DBLP:journals/tcad/LeeS98}.

\paragraph{Safety Theorems}
\label{sec:intro:contrib:semantics:thms}

We establish key semantic guarantees that validate the soundness of the language design:
%
\begin{itemize}
    \item \emph{Productivity \Cref{thm:semantics:grfrp:theorems:productivity}.} The step function of
          a service is total, it produces a valid output and new state for every possible state and
          input.
    \item \emph{Determinism \Cref{thm:semantics:grfrp:theorems:determinism}.} The output behavior of
          a service is uniquely determined by its input history.
    \item \emph{Equivalence \Cref{thm:semantics:grfrp:theorems:equivalence}.} The operational and
          denotational semantics are provably consistent: an infinite execution of the operational
          semantics over input streams produces outputs that satisfy the denotational semantics, and
          vice versa.
    \item \emph{Oversampling \Cref{thm:semantics:grfrp:theorems:oversampling}.} The outputs of a
          reactive service are unaffected by evaluations that occur when no new inputs have arrived,
          which justifies using an event-driven execution model instead of a periodic one without
          changing the behavior of the system.
\end{itemize}
%
These results have been proven by hand and are included in this manuscript. Unlike
\Velus~\cite{DBLP:conf/pldi/BourkeBDLPR17}, which formalizes correctness in \Rocq, our approach does not aim to
produce computer-formalized and verified proofs. Such formalization, while valuable, is known to be
labor-intensive and highly sensitive to implementation details, making it difficult to maintain
alongside evolving compiler code.

\subsection{Verification of \GRust Programs}
\label{sec:intro:contrib:greusot}

The \iso specification encourages the formal verification of \Asil system implementations. Typical
objectives include proving the absence of division-by-zero errors, avoiding overflows, and ensuring
critical boolean properties.

At \ampere, the current process is as follows: the system team models the system in \Simulink from
requirements; the safety team models the property in \textsc{Matlab} using the same requirements;
and the model checker \sldv verifies the property. \sldv also detects overflows, dead logic, and
possible division-by-zero errors.
%
These proofs apply to the high-level model, which aids comprehension, but the semantics of the
generated \Ci code are not verified. While full end-to-end proofs from source to binary are
unrealistic, \Ci does offer verification tools such as \textsc{Frama-C}~\cite{framac} that could be
used for code-level checks.

In contrast, the \Rust language benefits from the program verifier
\Creusot~\cite{DBLP:conf/icfem/DenisJM22}.
It leverages \Rust's compile-time memory safety checks (borrow checker) to simplify proof efforts.
For example, in \textsc{Frama-C}, the programmer must explicitly prove that pointers refer to
disjoint memory before writes, whereas in \Rust, the compiler enforces this by design, allowing
\Creusot proofs to focus solely on the program's logical correctness.

However, \Creusot operates at the level of plain \Rust code, making it non-trivial to express and
verify properties written for higher-level abstractions such as those in \GRust. To bridge this gap
and extend \GRust's applicability in the automotive domain, we developed \GReusot, a wrapper for
\Creusot.
%
\GReusot enables developers to state safety properties for \GRSync functions and components directly
in \GRust programs.
The prototype compiler automatically translates these specifications into \Rust-compatible
\Creusot annotations, integrating verification seamlessly into the development workflow.
%
This approach preserves the abstraction level of \GRust while ensuring that critical properties can
be formally checked on the generated code, contributing to \iso-compliant verification processes.

\section{Organization}
\label{sec:intro:orga}

This manuscript presents \GRust, from its practical design to its theoretical foundations. The
chapters are structured to guide the reader through the language specification, its formal
semantics, and a summary of the contributions.

\paragraph*{\Cref{cha:sota}} motivates the definition of \GRust by showing that the state of the art
on the domain does not completely address the problems emphasized above. It also details the
distinctive features of the language.

\paragraph*{\Cref{cha:grust}} provides a technical guide to the \GRust language, written for
practitioners and system developers. It begins with a practical example, an Adaptive Cruise Control
(ACC) system, to introduce the syntax and core concepts. The chapter then delves into the language
specification, detailing the grammar in EBNF, the type system, and identifier rules. It also covers
key aspects of the compiler implementation and the structure of the generated code. The chapter
concludes with a set of benchmarks that assess the practical efficiency of \GRust generated code,
comparing its performance against both state-of-the-art academic synchronous languages and
handwritten industrial \Rust code.

\paragraph*{\Cref{cha:semantics}} establishes the theoretical foundations of \GRust, and is aimed at
readers with an interest in formal methods and language design. We develop two complementary
semantic models: an operational semantics specifying the step-by-step execution of programs, and a
denotational semantics defining their mathematical meaning as streams that capture the history of
values over infinite executions. We formally prove the equivalence of these models and establish
key properties such as productivity and determinism. Furthermore, this chapter introduces a novel
oversampling theorem that guarantees the correctness of the execution strategy, a theoretical result
validated by the benchmarks in~\Cref{cha:grust}.

\paragraph*{\Cref{cha:conclusion}} concludes the manuscript by summarizing the contributions of this
thesis. It reflects on the current limitations of \GRust, and outlines promising directions for
future research. These include potential language extensions, enhanced verification capabilities,
and strategies for promoting industrial adoption.
